\section{Introduction}

The B method~\cite{Bmethod} and its associated toolsets, notably Atelier~B \cite{AtelierB}, have been used for decades in industrial developments where strong assurance arguments are required.
The B method is a formal development methodology based on stepwise refinement of abstract high-level specifications into executable code, supported by a rich mathematical language rooted in set theory and first-order logic.
B specifications are written as \emph{abstract machines} with sets, constants, variables, properties, invariants, and operations; Atelier~B generates \emph{proof obligations} (POs) whose discharge establishes, for example, invariant preservation and refinement correctness, and provides automated and interactive provers tailored to this logic.
In many certified developments, these tools are trusted as part of the verification chain.

In parallel, interactive theorem provers have seen rapid adoption, driven by their expressive type theories, powerful automation, and growing libraries of reusable mathematics and verification components.
Lean~4~\cite{lean} is based on the calculus of inductive constructions with quotients and has a dependent type system.
It offers a small trusted kernel, a standard library, and a metaprogramming framework that supports its use as both a programming language and a proof assistant.
However, despite B’s long-standing industrial role, support for using modern interactive theorem provers as backends for Atelier~B remains very limited: some previous work targets Isabelle~\cite{IsabelleHOL} but is either no longer maintained~\cite{Schmalz2012RodinIsabelle} or provides specific support~\cite{Wolff2024EB_Isabelle} for Event-B~\cite{AbrialEventB} in Rodin~\cite{AbrialRodin}.
Earlier work also studied a typed functional semantics of B and its translation to PVS~\cite{Bodeveix2002BtoPVS}.

This paper presents \emph{\barrel} (B Automated tRanslation for Reasoning in Lean), a Lean~4 library that bridges Atelier~B and Lean by turning B proof obligations into Lean theorem statements to be proved by automation or interactively.
Given a B machine (or an existing POG file), \barrel invokes Atelier~B to generate POs when needed and translates them into Lean theorems expressed over Mathlib~\cite{mathlib4}'s set-theoretic primitives, preserving the original structures and notations as closely as possible.
\barrel also generates its own \emph{well-definedness} (WD) proof obligations, which capture the definedness side conditions of B partial operators (e.g.,\ minimum, function application, and cardinality).
These are treated as standard Lean statements, so that definedness constraints can be proved once and then reused in subsequent interactive proofs.
The library provides commands for importing B developments and proving their obligations with standard Lean tactic scripts; the resulting theorems are added to the environment under names derived from the original POG tags.
The translation pipeline is organized into clearly separated components: a lightweight embedding of B syntax, a POG file reader, an encoding layer mapping B constructs to Lean expressions, and a discharger that generates theorem declarations and connects them to user proofs. This high-level architecture is illustrated in \cref{fig:barrel-archi}.



\begin{figure}[t]
	\centering
	\begin{tikzpicture}[
			file/.style={chamfered rectangle, draw, chamfered rectangle corners=north west},
			entity/.style={rectangle, draw},
			node distance=1.5cm and 1.5cm,
			every node/.style={inner sep=6pt, align=center, minimum height=1.22cm, outer sep=0},
			% start chain = going right,
			arr/.style={->, shorten <=3pt, shorten >=3pt}
		]
		\node[file] (B mch) {.mch};
		\node[file, right=of B mch] (B bxml) {.bxml};
		\node[entity, right=of B bxml] (pos) {Proof\\obligations};
		\node[entity, right=of pos] (thms) {Lean\\theorems};

		\draw[arr] (B mch) -- node[above=-4mm, midway] {\scriptsize Parsing} (B bxml);
		\draw[arr] (B bxml) -- node[above=-6mm, midway] {\scriptsize PO\\[-.6ex]\scriptsize generation} (pos);
		\draw[arr] (pos) -- node[above=-6mm, midway] {\scriptsize PO\\[-.6ex]\scriptsize encoding} (thms);
		\draw[arr, shorten <=7pt, shorten >=7pt] (thms.south) arc[start angle=220, end angle=500, x radius=9.2mm, y radius=9.5mm] node[midway, fill=white, minimum height=0pt, inner sep=2pt] {\scriptsize Proving};
	\end{tikzpicture}
	\caption{High-level picture of \barrel.}
	\label{fig:barrel-archi}
\end{figure}

Our contributions are threefold:
\begin{itemize}
	\item \barrel, an open-source, modular backend for B developments, written in Lean~4;
	\item a translation pipeline preserving B notations and proof obligations, with explicit well-definedness conditions for partial operators;
	\item an analysis of proof checking and modular reuse, combining two Atelier~B examples with a complete B refinement development.
\end{itemize}
Potential extensions include a verified proof obligation generator and an embedded DSL for writing B developments directly in Lean.
% We do not discuss semantics or soundness of the encoding in this paper.

The rest of this paper is organized as follows.
After reviewing related work, we recall the necessary background on the B method and present the design and implementation of \barrel, including its treatment of partial operators and basic automation.
We then examine the checks required during interactive proof (\cref{sec:proof-safety}), motivating the comparison of lemma reuse, refinement dependencies, and automation in the leader-election case study (\cref{sec:case-study}).
We conclude with future extensions and the remaining trust assumptions.
